% Volúmenes persistentes y variables
\subsection{Volúmenes persistentes y variables}\label{sec:5_2_volumenes_variables}

\subsubsection{Volumen de datos en /apps/data/postgresql/defensoria\_db}
Un contenedor se puede borrar y volver a crear en cualquier momento; si los datos vivieran dentro de él, se perderían. Por eso la carpeta donde PostgreSQL guarda la base (\texttt{/var/lib/postgresql/data} dentro del contenedor) se monta desde el disco de la VPS:
\begin{lstlisting}
-v /apps/data/postgresql/defensoria_db:/var/lib/postgresql/data:Z
\end{lstlisting}
\begin{itemize}
  \item La carpeta del host la crea el script la primera vez; al iniciar, PostgreSQL le asigna su propio usuario (\texttt{postgres}, UID 70), por eso no se modifica a mano (sección \ref{sec:2_3_directorios_permisos}).
  \item Detener, borrar o recrear el contenedor (\texttt{down} / \texttt{up-defensoria}) \textbf{conserva} todos los datos.
  \item Solo el comando \texttt{reset} del script borra esta carpeta, y pide escribir \texttt{BORRAR} para confirmarlo.
  \item La opción \texttt{:Z} ajusta la etiqueta de seguridad del volumen para sistemas con SELinux; en Ubuntu no tiene efecto, pero se conserva para que el script funcione igual en otras distribuciones.
\end{itemize}

\subsubsection{Script de inicialización 00-init.sql y scripts semilla}
La carpeta \nolinkurl{/apps/aplicaciones/defensoria/database-scripts} tiene dos subcarpetas con propósitos distintos:

\begin{table}[H]
\centering
\caption{Scripts SQL de la base de datos}
\small
\begin{tabular}{p{2.2cm}p{4.6cm}p{7.6cm}}
\toprule
\textbf{Carpeta} & \textbf{Archivos} & \textbf{Cuándo se ejecutan} \\
\midrule
\texttt{init/} & \texttt{00-init.sql} & Solo una vez, automáticamente, cuando el contenedor se crea con la carpeta de datos vacía. Se monta en \texttt{/docker-entrypoint-initdb.d} (solo lectura). \\
\texttt{seeds/} & \texttt{01-dependencias\_seed.sql}, \texttt{02-chatbot\_seed.sql}, \texttt{03-personal\_test.sql}, \texttt{04-recepcionistas.sql}, \texttt{05-abogados.sql} & A mano, después de levantar los microservicios, porque necesitan que las tablas ya existan (sección \ref{sec:6_5_datos_semilla}). \\
\bottomrule
\end{tabular}
\end{table}

\texttt{00-init.sql} solo configura la base; no crea tablas, porque las crea cada microservicio al arrancar (Hibernate con \texttt{ddl-auto: update}):
\begin{lstlisting}
CREATE EXTENSION IF NOT EXISTS "uuid-ossp";
ALTER DATABASE defensoria_db SET timezone TO 'America/Mexico_City';
\end{lstlisting}
\begin{itemize}
  \item \texttt{uuid-ossp}: extensión que genera identificadores únicos (UUID) que usan algunas tablas.
  \item \texttt{timezone}: todas las fechas que guarda la base quedan en hora del centro de México.
\end{itemize}
\begin{lstlisting}
# ls -l /apps/aplicaciones/defensoria/database-scripts/*
# cat /apps/aplicaciones/defensoria/database-scripts/init/00-init.sql
\end{lstlisting}
\captura{bd-02-scripts-sql}{Scripts de inicialización y semilla en el VPS Base de Datos}
